\documentclass[english,circlenumberstyle]{irtesbeamer}

\usepackage[utf8]{inputenc}
\usepackage{algpseudocode}
\usepackage{multicol}
\usepackage{tabularx}

\title{Documentation of the IRTES Theme for Beamer}
\subtitle{Macros provided by the theme}

\author[S.~Galland]{St\'ephane~Galland}

\finalslidetext{End of the Documentation}

\begin{document}

%%%%%%%%%%%%%%%%%%%%%%%%%%%%%%%%%%%%%%%%%%%%%%%%%
\tableofcontentslide

%%%%%%%%%%%%%%%%%%%%%%%%%%%%%%%%%%%%%%%%%%%%%%%%%
\section{Introduction}
\begin{frame}{\secname}
	\begin{itemize}
	\item This document explain the {\LaTeX} macros that are provided by the IRTES theme for beamer.
	\item To use the IRTES theme for beamer, you must:
		\begin{enumerate}
		\item Install the files of the IRTES theme inside your {\LaTeX} distribution.
		\item Create a document, which uses the document class: \texttt{irtesbeamer}
		\end{enumerate}
	\end{itemize}
\end{frame}

\begin{frame}{Major Features}
	\begin{itemize}
	\item Apply the official style for the IRTES institute.
	\item Add a title slide (with label \texttt{titleslide} for \texttt{{\textbackslash}hyperlink}) as the first slide.
	\item Add a final slide at the end of the presentation to avoid ``black'' screen at the end of the presentation.
	\end{itemize}
\end{frame}

\section{Document Attributes}

\subsection{Class Options}
\begin{frame}{\subsecname~(1/2)}\smaller
	The document class \texttt{irtesbeamer} supports the following options:
	\begin{itemize}
	\item \texttt{english}: the slides are written in English (this is the default).
	\item \texttt{french}: the slides are written in French.
	\item \texttt{nocover}: do not add a slide for the title and the ``thanks''.
	\item \texttt{thanksslide}: add a slide with the word ``Thanks'' at the end (this is the default). 
	\item \texttt{questionslide}: add a slide with the word ``Question'' at the end. 
	\item \texttt{repeattitleslide}: add the title slide at the end.
	\item \texttt{textnumberstyle}: the frame numbers are output as text.\hfill\hyperlink{progressbartypes}{\beamergotobutton{See example}}
	\item \texttt{circlenumberstyle}: the frame numbers are output in a circle.\hfill\hyperlink{progressbartypes}{\beamergotobutton{See example}}
	\item \texttt{sectioncirclenumberstyle}: the frame numbers are output in a circle.\hfill\hyperlink{progressbartypes}{\beamergotobutton{See example}}
 	\end{itemize}
\end{frame}

\begin{frame}{\subsecname~(2/2)}\smaller
	\begin{itemize}
	\item \texttt{partcirclenumberstyle}: similar to \texttt{circlenumberstyle} except that the indicator is reset at the start of each part.
	\item \texttt{partsectioncirclenumberstyle}: similar to \texttt{sectioncirclenumberstyle} except that the indicator is reset at the start of each part.
	\item \texttt{eventintitlebar}: on the title slide, the event's name should be put in the title bar.
	\item \texttt{eventbelowauthors}: on the title slide, the event's name should be put below the authors.
	\item All the options accepted by Beamer.
 	\end{itemize}
\end{frame}

\subsection{Title and Subtitle}
\begin{frame}{\subsecname}
	\begin{block}{Title}
		Use the macro: \texttt{{\textbackslash}title\{title of the document\}}
	\end{block}
	\begin{block}{Subitle}
		Use the macro: \texttt{{\textbackslash}subtitle\{subtitle of the document\}}
	\end{block}
\end{frame}

\subsection{Authors}
\begin{frame}{\subsecname}
	\begin{itemize}
	\item Authors are defined with the macro: \\
		\texttt{{\textbackslash}author[short]\{long\}}
		\begin{itemize}
		\item the \texttt{long} list of authors is displayed on the front page. You should separate the names with the macro \texttt{{\textbackslash}and}.
		\item the \texttt{short} list of authors is displayed inside the foot line of the slides. You \alert{must not separate} the names with the macro \texttt{{\textbackslash}and}.
		\end{itemize}
	\item Alternatively, you could define the authors with the macros: \\
		\texttt{{\textbackslash}addauthor[short]\{long\}} \\
		\texttt{{\textbackslash}addauthor*[short]\{long\}}
		\begin{itemize}
		\item Add \textbf{one} author to the list of the authors.
		\item The ``starred'' version applies some visual indicators to the name. It may be used to indicate the name of the talker for example.
		\end{itemize}
	\end{itemize}
\end{frame}

\subsection{Event Name}
\begin{frame}{\subsecname}
	\begin{block}{Definition of the event name}
		\begin{itemize}
		\item You could specify the name of the event for which the slides are written: \\
			\texttt{{\textbackslash}event[short]\{full\}}
		\item the \texttt{full} name is displayed on the front page.
		\item the \texttt{short} name is displayed inside the foot line of the slides.
		\end{itemize}
	\end{block}
	\begin{block}{Insert the event name in your slides}
		\begin{itemize}
		\item You could insert the full name of the event with: \\
			\texttt{{\textbackslash}inserteventname}
		\item You could insert the short name of the event with: \\
			\texttt{{\textbackslash}insertshorteventname}
		\end{itemize}
	\end{block}
\end{frame}

\subsection{Institute Name and URL}

\begin{frame}{Institute Name}
	\begin{itemize}
	\item You could change the name of the institute with: \\
		\texttt{{\textbackslash}institute[short]\{full\}}
	\item the \texttt{full} name is displayed on the front page.
	\item the \texttt{short} name is displayed inside the foot line of the slides.
	\item The default full name is: \\
		\textbf{Universit\'e de Technologie de Belfort-Montb\'eliard}, 90010 Belfort cedex, France - \url{http://irtes.utbm.fr}
	\item The default short name is: IRTES.
	\end{itemize}
\end{frame}

\begin{frame}{Institute URL}
	\begin{itemize}
	\item You could change the URL of the institute with: \\
		\texttt{{\textbackslash}instituteurl\{url\}}
	\item The default url is: \url{http://irtes.utbm.fr}
	\item To insert the institute's URL, you could use the macro: \\
			\texttt{{\textbackslash}insertinstituteurl}
	\end{itemize}
\end{frame}

\subsection{Partner Logo}
\begin{frame}{\subsecname}
	\begin{itemize}
	\item You could add on all slides the logo of your partner: \\
		\texttt{{\textbackslash}partnerlogo[options]\{filename\}}
	\item The IRTES theme has enough space for only one logo with a width of 2cm.
	\item The \texttt{filename} is the name of the picture.
	\item This figure is declared with \texttt{{\textbackslash}pgfdeclareimage} with the key ``IRTESpartnerlogo''.
	\item The \texttt{options} are passed to \texttt{{\textbackslash}pgfdeclareimage}. The default option is \texttt{width=2cm}.
	\end{itemize}
\end{frame}

\section{Document Structure}

\subsection{Part Sectioning}
\begin{frame}{\subsecname}
	\begin{itemize}
	\item The IRTES theme provides a specific implementation of the \texttt{{\textbackslash}part} macro: \\
		\texttt{{\textbackslash}part[options]{title}} \\
		\texttt{{\textbackslash}part*[options]{title}} \\
	\item Options may be: \begin{itemize}
		\item a string value that is the ``short'' title of the part that is appearing in the table of contents.
		\item the pair \texttt{title=text} to define the ``short'' title.
		\item the pair \texttt{author=text} to define the author of the part.
		\end{itemize}
	\item The stared version of \texttt{{\textbackslash}part} does not add the part in the table of contents.
	\end{itemize}
\end{frame}

\subsection{Table of Contents}
\begin{frame}{\subsecname}
	\begin{itemize}
	\item The IRTES theme provides a convenient macro to insert a table of contents: \\
		\texttt{{\textbackslash}tableofcontentsslide[options]}
	\item It is equivalent to: \\
		\texttt{{\textbackslash}begin\{frame\}[c]} \\
		\texttt{{\textbackslash}frametitle\{{\textbackslash}translate\{Outline\}\}} \\
		\texttt{{\textbackslash}tableofcontents[options]} \\
		\texttt{{\textbackslash}end\{frame\}}
	\item The option \texttt{onlyparts} permits to display the list of the parts, only.
	\end{itemize}
\end{frame}

\subsection{Picture Slide}
\begin{frame}{\subsecname}
	The IRTES theme provides a macr othat permits to display a picture on the entire slide:
	\begin{itemize}
	\item \texttt{{\textbackslash}figureslide[width=XXX,height=YYY]\{Title of the slide\}\{file\}}
	\item Example: \texttt{{\textbackslash}figureslide\{XyZ\}\{irteslogo\}} 
	\end{itemize}
	\vspace{2cm}
	\begin{center}
		\begin{minipage}{.25\paperwidth}
			\begin{picture}(0,0)
				\put(0,0){\includegraphics[width=\linewidth]{IRTESbackground}}
				\put(22,59){\textcolor{white}{\tiny XyZ}}
				\put(5,20){\includegraphics[width=.9\linewidth]{irteslogo}}
			\end{picture}
		\end{minipage}
	\end{center}
\end{frame}

\subsection{Final Slide}
\begin{frame}{A Conclusion Slide}
	\begin{itemize}
	\item The IRTES theme automatically adds a slide at the end of the presentation to avoid ``black'' screen.
	\item The default text on this slide is: ``\translate{Thanks}''.
	\item An other text that is available is: ``\translate{Questions}''.
	\item The class options \texttt{thanksslide} and \texttt{questionslide} permits to select one of these possibilities.
	\item The macro \texttt{{\textbackslash}finalslidetext\{text\}} may be used to set the text by hand.
	\item \emph{This slide has the label \texttt{finalslide} for \texttt{{\textbackslash}hyperlink}.}
	\end{itemize}
\end{frame}

\subsection{Appendix}
\begin{frame}{\subsecname}
	\begin{itemize}
	\item The IRTES theme supports the appendix.
	\item To create the appendix, you must:
		\begin{enumerate}
		\item put the macro \texttt{{\textbackslash}appendix} in your {\TeX} file; or
		\item put the macro \texttt{{\textbackslash}bibliography} in your {\TeX} file.
		\end{enumerate}
	\item All the slides put after the creation of the appendix by one of these methods are assumed to be part of the appendix.
	\item The slides in the appendix are not considered in the total number of slides for the core part of the document (see \texttt{{\textbackslash}inserttotalcoreframenumber}).
	\item The slide numbers in the appendix are roman (not arabic), and the numbers are reset.
	\end{itemize}
\end{frame}

\subsection{Bibliography and Library}
\begin{frame}{Bibliography}
	\begin{itemize}
	\item You are able to include a bibliography in your slides with the two standard {\LaTeX} macros: \\
		\texttt{{\textbackslash}bibliographystyle\{style\}} \\
		\texttt{{\textbackslash}bibliography\{filename\}}
	\item The default style is \texttt{apalike}.
	\item When the macro \texttt{{\textbackslash}bibliography} is used, the appendix section is started if it was not already done.
	\item The bibliography slides are \alert{always} at the end of the document. Even if you put slides after the {\textbackslash}bibliography macro.
	\item \emph{The first slide of the bibliography has the label \texttt{bibliographyslide} for \texttt{{\textbackslash}hyperlink}.}
	\end{itemize}
\end{frame}

\begin{frame}{Library}
	\begin{itemize}
	\item You are able to include a list of books in your presentation with the macro: \\
		\texttt{{\textbackslash}libraryslide[subtitle or comment]\{picture\}} \\
		\texttt{\{title\}\{authors\}\{How published\}\{ISBN\}}
	\item The macro create a slide for a book.
	\end{itemize}
\end{frame}

\section{Slide Content}

\subsection{Colors}
\begin{frame}{\subsecname}
	The IRTES theme defines the following colors:
	\begin{itemize}
	\item \texttt{IRTESblue} : \textcolor{IRTESblue}{Example}
	\item \texttt{IRTESmagenta} : \textcolor{IRTESmagenta}{Example}
	\item \texttt{IRTESlightblue} : \textcolor{IRTESlightblue}{Example}
	\item \texttt{IRTESdarkgray} : \textcolor{IRTESdarkgray}{Example}
	\end{itemize}
\end{frame}

\subsection{Alert Boxes}
\begin{frame}{\subsecname}
	The IRTES theme defines the a box for alerts: \\
	\begin{itemize}
	\item \texttt{{\textbackslash}alertbox\ensuremath{<}frame\_spec\ensuremath{>}\{this is an alert text\}} \\[.5cm]
		\alertbox{this is an alert text}
		\vspace{1cm}
	\item \texttt{{\textbackslash}alertbox*\ensuremath{<}frame\_spec\ensuremath{>}\{this is an alert text\}} \\[.5cm]
		\alertbox*{this is an alert text}
	\item Note that \texttt{\ensuremath{<}frame\_spec\ensuremath{>}} is optional. It permits to specify the Beamer frame in which the box is displayed.
	\end{itemize}
\end{frame}

\subsection{Enhanced Description}
\begin{frame}{\subsecname}
	The IRTES theme provides an enhanced definition of the \texttt{description} environment. \\[.5cm]
	{\smaller
	\texttt{{\textbackslash}begin\{description\}} \\
	\texttt{{\textbackslash}item text1} \\
	\texttt{{\textbackslash}item[Item Name] text2} \\
	\texttt{{\textbackslash}item$<$frame\_spec$>$ text1} \\
	\texttt{{\textbackslash}item$<$frame\_spec$>$[Item Name] text2} \\
	\texttt{{\textbackslash}end\{description\}}} \\[.5cm]
	\begin{description}
	\item text1
	\item[Item Name] text2
	\item<2> text1
	\item<2>[Item Name] text2
	\end{description}
\end{frame}

\subsection{Enhanced Enumeration}
\begin{frame}{\subsecname (1/2)}
	The IRTES theme provides an enhanced definition of the \texttt{enumerate} environment. \\[.5cm]
	{\smaller
	\texttt{{\textbackslash}begin\{enumerate\}[counter format]} \\
	\texttt{{\textbackslash}item text1} \\
	\texttt{{\textbackslash}item[Item Name] text2} \\
	\texttt{{\textbackslash}item$<$frame\_spec$>$ text1} \\
	\texttt{{\textbackslash}item$<$frame\_spec$>$[Item Name] text2} \\
	\texttt{{\textbackslash}end\{enumerate\}}} \\[.5cm]
	\begin{enumerate}
	\item text1
	\item[Item Name] text2
	\item<2> text1
	\item<2>[Item Name] text2
	\end{enumerate}
\end{frame}

\begin{frame}{\subsecname (2/2)}
	Below, the optional parameter \texttt{counter format} is set to \texttt{"a)"}: \\[.5cm]
	\begin{enumerate}[a)]
	\item text1
	\item[Item Name] text2
	\item<2> text1
	\item<2>[Item Name] text2
	\end{enumerate}
\end{frame}

\subsection{Text Size}

\begin{frame}{\subsecname}
	In addition to the standard Beamer macros for selected the text side (\texttt{{\textbackslash}TINY}, \texttt{{\textbackslash}Tiny}, \texttt{{\textbackslash}tiny}, \texttt{{\textbackslash}footnotesize}, \texttt{{\textbackslash}scriptisize}, \texttt{{\textbackslash}small}, \texttt{{\textbackslash}normalsize}, \texttt{{\textbackslash}large}, \texttt{{\textbackslash}Large}, \texttt{{\textbackslash}huge}, \texttt{{\textbackslash}Huge}), the IRTEs theme enables the following macros:
	\begin{itemize}
	\item \texttt{{\textbackslash}smaller} : to decrease the size of the text, and
	\item \texttt{{\textbackslash}larger} : to increase the size of the text.
	\end{itemize}
\end{frame}
\makeatother

\subsection{Underlining}

\makeatletter
\begin{frame}{\subsecname}
	\begin{itemize}
	\item The IRTES theme redefines the macro \texttt{{\textbackslash}underline} to move the line closer to the text.
	\item Before: \beamer@theme@irtes@oldunderline{Example}
	\item After: \underline{Example}
	\end{itemize}
\end{frame}
\makeatother

\subsection{Additional Symbols}
\begin{frame}{\subsecname}
	\begin{itemize}
	\item The IRTES theme (re)defines the macros for several symbols
	\end{itemize}
	\begin{tabularx}{\linewidth}{|l|X|l|X|}
	\hline
	\multicolumn{2}{|c|}{\textbf{From IRTES theme}} & \multicolumn{2}{c|}{\textbf{From \TeX}} \\
	\hline
	\texttt{{\textbackslash}copyright} & \copyright & \texttt{{\textbackslash}copyright} & \textcopyright \\
	\hline
	\texttt{{\textbackslash}trademark} & \trademark & \texttt{{\textbackslash}texttrademark} & \texttrademark \\
	\hline
	\texttt{{\textbackslash}regmark} & \regmark & \texttt{{\textbackslash}textregistered} & \textregistered \\
	\hline
	\end{tabularx}
\end{frame}

\subsection{Exponents and Indices}
\begin{frame}{\subsecname}
	\begin{itemize}
	\item The IRTES theme (re)defines the macros to put text in exponent or in indice.
	\item The macros \texttt{{\textbackslash}textup} and \texttt{{\textbackslash}textdown} try to add a space after the text, when it is allowed by the typographic rules (it uses the macro \texttt{{\textbackslash}xspace}).
	\end{itemize}
	\begin{tabularx}{\linewidth}{|l|X|l|X|}
	\hline
	\multicolumn{2}{|c|}{\textbf{From IRTES theme}} & \multicolumn{2}{c|}{\textbf{From \TeX}} \\
	\hline
	\texttt{{\textbackslash}textup} & ABC\textup{ABC}D & \texttt{{\textbackslash}textsuperscript} & ABC\textsuperscript{ABC}D \\
	\hline
	\texttt{{\textbackslash}textsubscript} & ABC\textsubscript{ABC}D & - & - \\
	\hline
	\texttt{{\textbackslash}textdown} & ABC\textdown{ABC}D & - & - \\
	\hline
	\end{tabularx}
\end{frame}

\subsection{Footnotes}
\begin{frame}{\subsecname}
	\begin{itemize}
	\item The IRTES theme provides two versions of the footnote macro:
		\begin{itemize}
		\item \texttt{{\textbackslash}footnote\{text\}} shows a footnote with a number, and
		\item \texttt{{\textbackslash}footnote*\{text\}} shows a footnote without a number.
		\end{itemize}
	\item Additionnally, a footnote with bibliography citation may be added:
		\begin{itemize}
		\item \texttt{{\textbackslash}footcite\{keys\}} shows the given citations in a footnote.
		\end{itemize}
	\end{itemize}
\end{frame}

\subsection{Side Notes}
\begin{frame}{\subsecname}
	\begin{itemize}
	\item The IRTES theme provides a macro to put a text on the side of the frame \\
		\texttt{{\textbackslash}sidenode\{text\}}
	\item Example: \texttt{{\textbackslash}sidenote\{text on the side\}}
		\sidenote{text on the side}
	\end{itemize}
\end{frame}

\subsection{Localized Quotes}
\begin{frame}{\subsecname}
	\begin{itemize}
	\item The IRTES theme provides macros to output localized quotes:
		\begin{description}
		\item[English] \texttt{{\textbackslash}ukquote{text}} \\
			Example: \ukquote{text}
		\item[French] \texttt{{\textbackslash}frquote{text}} \\
			Example: \frquote{text}
		\item[Latin] \texttt{{\textbackslash}latquote{text}} \\
			Example: \latquote{text}
		\end{description}
	\item The following macros are used by the quote macros:
		\begin{itemize}
		\item \texttt{{\textbackslash}textgravedbl} : \textgravedbl
		\item \texttt{{\textbackslash}textacutedbl} : \textacutedbl
		\item \texttt{{\textbackslash}guillemotleft} : \guillemotleft
		\item \texttt{{\textbackslash}guillemotright} : \guillemotright
		\end{itemize}
	\end{itemize}
\end{frame}

\subsection{Frame Number}
\begin{frame}[label=progressbartypes]{\subsecname}\tiny
	\begin{itemize}
	\item The total number of slides in the core part of the presentation could be obtained with: \\
		\texttt{{\textbackslash}inserttotalcoreframenumber}
	\item The following macro provides the progression in the slides:
		\begin{itemize}
		\item \texttt{{\textbackslash}insertframenumbering[type number]}
		\item The \texttt{type number} is the identifier of the progress bar to be used:
			\begin{tabularx}{\linewidth}{|c|c|X|}
			\hline
			\textbf{Type number} & \textbf{Output} & \textbf{Explanation} \\
			\hline
			\texttt{0} & \insertframenumbering[0] & Show the current frame and its position in the total number of frames. \\
			\hline
			\texttt{1} & \colorbox{IRTESlightblue}{\insertframenumbering[1]} & Show the current frame and the total number of frames. \\
			\hline
			\texttt{2} & \insertframenumbering[2] & Same as the type \texttt{0} with a progression bar for the current section. \\
			\hline
			\end{tabularx}
		\end{itemize}
	\end{itemize}
\end{frame}

\subsection{Multimedia Link}
\begin{frame}{\subsecname}
	\begin{itemize}
	\item The IRTES theme provides a convenient macro to include links to multimedia resources.
	\item \texttt{{\textbackslash}videolink[options]\{resource\_path\}\{img\_path\}}
		\begin{itemize}
		\item Display a picture with a link button. When the user click on the picture, the resource is run (viewed).
		\item \texttt{options} are the options to pass to the \texttt{{\textbackslash}includegraphics} (width...)
		\item \texttt{resource\_path} is the path to the multimedia resource.
		\item \texttt{img\_path} is the path to the picture to display in the slide.
		\end{itemize}
	\end{itemize}
	\begin{center}
		\videolink[width=.3\linewidth]{./documentation.tex}{IRTEStitle}
	\end{center}
\end{frame}


\subsection{Link to a Slide from an Picture}
\begin{frame}{\subsecname}
	\begin{itemize}
	\item The IRTES theme provides a convenient macro to create links to other slides with a picture in the link.
	\item \texttt{{\textbackslash}picturegoto[options]\{label\}\{img\_path\}}
		\begin{itemize}
		\item Display a picture with a link button. When the user click on the picture, the slide with the given label is displayed.
		\item \texttt{options} are the options to pass to the \texttt{{\textbackslash}includegraphics} (width...)
		\item \texttt{label} is the label of the target slide.
		\item \texttt{img\_path} is the path to the picture to display in the slide.
		\end{itemize}
	\end{itemize}
	\begin{center}
		\picturegoto[width=.3\linewidth]{titleslide}{IRTEStitle}
	\end{center}
\end{frame}

\subsection{Zooming on a \TeX\ Box}
\begin{frame}{\subsecname\ (1/2)}
	\begin{itemize}
	\item Beamer provides the macro \texttt{{\textbackslash}framezoom} to zoom on a part of a frame. \emph{But, it create a new slide and it is difficult to return to the original slide with a single click}.	
	\item The IRTES theme provides a new macro for zooming on click.
	\item \texttt{{\textbackslash}zoombox[options]\{box content\}}
		\begin{itemize}
		\item Display the content of the box. When clicked, display the box content after fitting it to the entire screen. When clicked again, the original slide is restore.
		\item \texttt{options} are:
			\begin{enumerate}
			\item \texttt{border=XXpt}: specify the size of the border lines around the box.
			\item \texttt{left=XXpt}: specify the size of the left margin.
			\item \texttt{right=XXpt}: specify the size of the right margin.
			\item \texttt{top=XXpt}: specify the size of the top margin.
			\item \texttt{bottom=XXpt}: specify the size of the bottom margin.
			\item \texttt{margin=XXpt}: specify the size of all of the margins.
			\end{enumerate}
		\end{itemize}
	\end{itemize}
\end{frame}

\begin{frame}{\subsecname\ (2/2)}
	\begin{alertblock}{Caution}
	The macro \texttt{{\textbackslash}zoombox} is working in viewers that are supporting JavaScript (Acrobat Reader...)
	\end{alertblock}
	\vfill
	\begin{center}
		\zoombox{ZOOMING EXAMPLE}
	\end{center}
\end{frame}

\subsection{Adjusting a Box}
\begin{frame}{\subsecname}
	\begin{itemize}
	\item The IRTES style provides the macro \texttt{{\textbackslash}adjustbox} to add margins to a box.
	\item \texttt{{\textbackslash}adjustbox[options]\{box content\}}
	\item The options are:
		\begin{itemize}
		\item \texttt{left=XXpt} is the size of the left margin.
		\item \texttt{right=XXpt} is the size of the right margin.
		\item \texttt{top=XXpt} is the size of the top margin.
		\item \texttt{bottom=XXpt} is the size of the bottom margin.
		\item \texttt{size=XXpt} is the size of all of the margins.
		\end{itemize}
	\end{itemize}
\end{frame}

\subsection{Put at absolute position}
\begin{frame}{\subsecname}
	\begin{itemize}
	\item If you want to pout something at an absolute position in your frame, you could use: \\
		\texttt{{\textbackslash}putat(x,y)\{something\}}
	\item Example: \\
		\texttt{{\textbackslash}putat(0,-150)\{{\textbackslash}color\{red\}\{TESTING\}\}}
	\end{itemize}
	\putat(0,-150){\color{red}{TESTING}}
\end{frame}

\end{document}
